% Source printed pp. 83-86.
\section{预层与层}\label{src:3.3}
\begin{epigraph}
按照 A. Grothendieck 的观点, 做几何其实并不需要空间, 只需要这个设想中的空间上的层范畴.
\epigraphauthor{Vladimir Berkovich~\cite{Ber90}}
\end{epigraph}
回顾, 若 $X$ 是拓扑空间, 则 $X$ 上的集合值预层就是开集范畴 $\Op(X)$ 上的反变函子 $F:\Op(X)\to\Sets$. 层公理可以简洁地表述为: 对每个覆盖 $U=\bigcup_i U_i$, 序列
\[
F(U)\longrightarrow\prod_i F(U_i)\rightrightarrows\prod_{i,j}F(U_i\cap U_j)
\]
正合 (即是等化子图), 其中两个映射 $F(U_i)\rightrightarrows F(U_i\cap U_j)$ 由沿两个包含映射 $U_i\cap U_j\subset U_i$ 及 $U_i\cap U_j\subseteq U_j$ 的限制诱导. 还要注意, 交集 $U_i\cap U_j$ 也可以看作纤维积 $U_i\times_X U_j$.

一般景上预层与层的定义仿照拓扑空间的情形. 利用下降, 我们将证明大平展景上一些最初的层的例子.
\subsection{定义}\label{src:3.3.1-subsection}
\begin{definition}[预层]\label{src:3.3.1}
范畴 $\mathcal S$ 上的\newterm{预层}{presheaf}是反变函子 $\mathcal S\to\Sets$.
\end{definition}
\begin{remark}\label{src:3.3.2}
若 $F:\mathcal S\to\Sets$ 是预层, 而 $f:S\to T$ 是 $\mathcal S$ 中的映射, 则元素 $b\in F(T)$ 的拉回 $F(f)(b)$ 有时记为 $f^*b$ 或 $b|_S$.
\end{remark}
\begin{example}[可表预层]\label{src:3.3.3}
若 $\mathcal S$ 是范畴且 $S\in\mathcal S$, 则 $\Mor_{\mathcal S}(-,S):\mathcal S\to\Sets$ 定义一个\newterm{可表预层}{representable presheaf}.
\end{example}
\begin{definition}[层]\label{src:3.3.4}
景 $\mathcal S$ 上的\newterm{层}{sheaf}是预层 $F:\mathcal S\to\Sets$, 使得对每个对象 $S\in\mathcal S$ 及覆盖 $\{S_i\to S\}\in\Cov(S)$, 序列
\begin{equation}\label{src:3.3.5}
F(S)\longrightarrow\prod_i F(S_i)\rightrightarrows\prod_{i,j}F(S_i\times_S S_j)
\end{equation}
正合, 其中两个映射 $F(S_i)\rightrightarrows F(S_i\times_S S_j)$ 由两个映射 $S_i\times_S S_j\to S_i$ 及 $S_i\times_S S_j\to S_i$ 诱导.
\end{definition}
预层或层的态射按定义是自然变换.
\begin{remark}\label{src:3.3.6}
\bookeqref{3.3.5}~的正合性是指它是等化子图: $F(S)$ 被等同于 $\prod_i F(S_i)$ 的一个子集, 由在两个映射 $F(S_i)\rightrightarrows F(S_i\times_S S_j)$ 下的像相等的元素组成.
\end{remark}
\begin{exercise}\label{src:3.3.7}
设 $F$ 是 $\Sch$ 上的预层. 证明下列条件等价:
\begin{enumerate}[(1)]
\item $F$ 是 $\Sch_{\text{ét}}$ (或 $\Sch_{\mathrm{fppf}}$, $\Sch_{\mathrm{fpqc}}$) 上的层.
\item $F$ 将余积变为积 (即对概形 $U_i$, 有 $F(\coprod_i U_i)=\prod_i F(U_i)$), 且对每个概形的满射平展 (或 fppf, 忠实平坦) 态射 $S'\to S$, 序列 $F(S)\to F(S')\rightrightarrows F(S'\times_S S')$ 正合.
\item $F$ 是大 Zariski 拓扑 $\Sch_{\mathrm{Zar}}$ 中的层, 且对每个仿射概形的满射平展 (或 fppf, 忠实平坦) 态射 $S'\to S$, 序列 $F(S)\to F(S')\rightrightarrows F(S'\times_S S')$ 正合.
\end{enumerate}
\textit{提示:} 对覆盖 $\{S_i\to S\}$, 考虑映射 $\coprod_i S_i\to S$.
\end{exercise}
\begin{proposition}[概形是层]\label{src:3.3.8}
若 $X\to S$ 是概形态射, 则 $\Mor_S(-,X):\Sch/S\to\Sets$ 是 $(\Sch/S)_{\mathrm{fpqc}}$ 上的层, 因而也是 $(\Sch/S)_{\text{ét}}$ 与 $(\Sch/S)_{\mathrm{fppf}}$ 上的层.
\end{proposition}
\begin{proof}
由于 $\Mor_S(-,X)$ 是大 Zariski 拓扑中的层, 由习题~\bookref{3.3.7}, 只需证明: 若 $T'\to T$ 是 $S$ 上仿射概形的忠实平坦态射, 则序列
\[
\Mor_S(T,X)\longrightarrow\Mor_S(T',X)\rightrightarrows\Mor_S(T'\times_T T',X)
\]
正合. 这恰是态射的 fpqc 下降 (推论~\bookref{3.1.10}).
\end{proof}
\begin{exercise}\label{src:3.3.9}
\begin{enumerate}[(a)]
\item 若 $F$ 与 $G$ 是景 $\mathcal S$ 上的层, 证明由 $S\mapsto\Mor_{\Sets}(F(S),G(S))$ 定义的预层 $\Mor_{\mathcal S}(F,G)$ 是 $\mathcal S$ 上的层.
\item 推出, 若 $X$ 与 $Y$ 是 $S$ 上的概形, 则将 $S$-概形 $T$ 送到 $\Mor_T(X_T,Y_T)$ 的函子 $\Mor_S(X,Y):\Sch/S\to\Sets$ 是 fpqc 拓扑中的层. (这给出推论~\bookref{3.1.10}(2) 的另一个证明.)
\end{enumerate}
\end{exercise}
\begin{exercise}[层的粘合]\label{src:3.3.10}
设 $\mathcal S$ 是景, $(X_i\to X)$ 是 $\mathcal S$ 中的覆盖. 若 $F_i$ 是限制景 $\mathcal S/X_i$ 上的层, 而 $\alpha_{ij}:F_i|_{X_{ij}}\to F_j|_{X_{ij}}$ 是 $\mathcal S/X_{ij}$ 上的层同构 (其中 $X_{ij}:=X_i\times_X X_j$), 在 $\mathcal S/X_{ijk}$ 上满足上循环条件 $\alpha_{jk}\circ\alpha_{ij}=\alpha_{ik}$ (其中 $X_{ijk}=X_i\times_X X_j\times_X X_k$), 证明存在唯一的 $\mathcal S/X$ 上的层 $F$, 以及满足 $\varphi_j|_{X_{ij}}=\alpha_{ij}\circ\varphi_i|_{X_{ij}}$ 的同构 $\varphi_i:F|_{X_i}\to F_i$.
\end{exercise}
\begin{exercise}\label{src:3.3.11}
证明, 概形的满射光滑 (或 fppf, fpqc) 态射是 $\Sch_{\text{ét}}$ (或 $\Sch_{\mathrm{fppf}}$, $\Sch_{\mathrm{fpqc}}$) 上层的满态射.
\end{exercise}

\subsection{纤维积}\label{src:3.3.2-subsection}
由 Yoneda 引理~(\bookref{1.3.12}), 若 $X\in\mathcal S$ 是范畴 $\mathcal S$ 的对象, 且 $F$ 是 $\mathcal S$ 上的预层, 则态射 $\alpha:X\to F$ (将其理解为预层态射 $\Mor(-,X)\to F$) 对应于 $F(X)$ 中的元素. 略微滥用记号, 我们也将这个元素记为 $\alpha$.
\begin{definition}\label{src:3.3.12}
给定范畴 $\mathcal S$ 上预层的态射 $\alpha:F\to G$ 与 $\beta:G'\to G$, $\alpha$ 与 $\beta$ 的纤维积是预层 $F\times_G G'$, 其在 $S\in\mathcal S$ 上的截面集合为 $F(S)\times_{G(S)}G'(S)$, 即
\begin{equation}\label{src:3.3.13}
\begin{aligned}
F\times_G G':\mathcal S&\longrightarrow\Sets,\\
S&\longmapsto\{(a,b)\in F(S)\times G'(S)\mid\alpha_S(a)=\beta_S(b)\}.
\end{aligned}
\end{equation}
\end{definition}
\begin{exercise}\label{src:3.3.14}
\begin{enumerate}[(a)]
\item 证明~\bookeqref{3.3.13}~是 $\Pre(\mathcal S)$ 中的纤维积 $F\times_G G'$. (这是习题~\bookref{1.3.41}~的推广, 但同样的证明应当适用.)
\item 证明, 若 $F,G,G'$ 都是景 $\mathcal S$ 上的层, 则 $F\times_G G'$ 也是. 特别地, \bookeqref{3.3.13}~也是 $\Sh(\mathcal S)$ 中的纤维积 $F\times_G G'$.
\end{enumerate}
\end{exercise}

\subsection{层化}\label{src:3.3.3-subsection}
\begin{theorem}[层化]\label{src:3.3.15}
设 $\mathcal S$ 是景. 遗忘函子 $\Sh(\mathcal S)\to\Pre(\mathcal S)$ 有左伴随 $F\mapsto F^{\mathrm{sh}}$, 称为\newterm{层化}{sheafification}.
\end{theorem}
\begin{proof}
$\mathcal S$ 上的预层 $F$ 称为\newterm{分离的}{separated}, 若对对象 $S$ 的每个覆盖 $\{S_i\to S\}$, 映射 $F(S)\to\prod_i F(S_i)$ 都是单射 (即截面若能粘合, 则粘合唯一). 记预层范畴与层范畴为 $\Pre(\mathcal S)$ 与 $\Sh(\mathcal S)$, 而 $\Pre_{\mathrm{sep}}(\mathcal S)\subset\Pre(\mathcal S)$ 是由分离预层组成的全子范畴. 我们将构造左伴随
\[
\begin{tikzcd}[column sep=huge]
\Sh(\mathcal S)\arrow[r,hook]&
\Pre_{\mathrm{sep}}(\mathcal S)\arrow[l,bend right=35,"\mathrm{sh}_2"']
\arrow[r,hook]&
\Pre(\mathcal S)\arrow[l,bend right=35,"\mathrm{sh}_1"']
\end{tikzcd}
\]
对 $F\in\Pre(\mathcal S)$, 定义 $\mathrm{sh}_1(F)$ 为 $S\mapsto F(S)/{\sim}$, 其中 $a\sim b$ 是指存在覆盖 $\{S_i\to S\}$, 使得对所有 $i$ 都有 $a|_{S_i}=b|_{S_i}$. 对 $F\in\Pre_{\mathrm{sep}}(\mathcal S)$, 定义 $\mathrm{sh}_2(F)$ 为
\[
S\longmapsto
\left\{\bigl(\{S_i\to S\},\{a_i\}\bigr)\ \middle|\
\begin{gathered}
\{S_i\to S\}\in\Cov(S),\quad a_i\in F(S_i),\\
a_i|_{S_{ij}}=a_j|_{S_{ij}}\text{ 对所有 }i,j\text{ 成立}
\end{gathered}
\right\}/{\sim},
\]
其中 $(\{S_i\to S\},\{a_i\})\sim(\{S_j'\to S\},\{a_j'\})$ 是指对所有 $i,j$ 都有 $a_i|_{S_i\times_S S_j'}=a_j'|_{S_i\times_S S_j'}$. 细节留给读者.
\end{proof}
\begin{remark}[拓扑斯]\label{src:3.3.16}
\newterm{拓扑斯}{topos}是等价于某个景上层范畴的范畴. 两个不同的景可能有等价的层范畴, 而拓扑斯是更基本的不变量. 拓扑斯在模理论中无疑很重要, 但在本讲义中不会起作用.
\end{remark}

\subsection{层为概形的判别法}\label{src:3.3.4-subsection}
下面是概形的 fppf 下降~(\bookref{3.1.11}, \bookref{3.1.12}~与~\bookref{3.1.14}) 的另一种表述.
\begin{proposition}[fppf 层为概形的下降判别法]\label{src:3.3.17}
设 $P$ 是概形态射的下列性质之一: 开浸入, 闭浸入, 仿射, 拟仿射, 或局部拟有限且分离. 设 $X\to Y$ 是概形的满射光滑 (或 fppf) 态射, $F$ 是 $(\Sch/Y)_{\text{ét}}$ (或 $(\Sch/Y)_{\mathrm{fppf}}$) 上的层. 考虑层的纤维积
\[
\begin{tikzcd}
F_X\arrow[r]\arrow[d]\arrow[dr,phantom,"\square"]&X\arrow[d]\\
F\arrow[r]&Y
\end{tikzcd}
\]
若 $F_X$ 是概形且 $F_X\to X$ 具有 $P$, 则 $F$ 是概形且 $F\to Y$ 具有 $P$.
\end{proposition}
\begin{proof}
由于 $F_X$ 是 $F$ 的拉回, $X\times_Y X$ 上有典范同构 $\alpha:p_1^*F_X\to p_2^*F_X$, 满足上循环条件 $p_{23}^*\alpha\circ p_{12}^*\alpha=p_{13}^*\alpha$. 由概形的 fppf 下降~(\bookref{3.1.11}, \bookref{3.1.12}~与~\bookref{3.1.14}), 存在具有 $P$ 的概形态射 $W\to Y$, 其拉回为 $F_X\to X$.
层 $F$ 可以等同于 $\Sch/Y$ 上如下定义的预层的层化:
\[
F^{\mathrm{pre}}:(T\to Y)\longmapsto
\operatorname{Eq}\bigl(F_X(T\to Y)\rightrightarrows
(F_X\times_X F_X)(T\to Y)\bigr).
\]
由概形是层~(\bookref{3.3.8}), $W$ 是层, 且有预层态射 $F^{\mathrm{pre}}\to W$. 利用层化的泛性质, 得到层态射 $\alpha:F\to W$, 它沿 $X\to Y$ 的拉回是同构 $\alpha_X:F_X\to W_X$.
不难看出, 这意味着 $\alpha$ 是同构. 例如, $\Hom(W,F)$ 是层 (习题~\bookref{3.3.9}), 而 $\alpha_X$ 的逆 $\beta_X$ 定义了 $X\to Y$ 上的截面, 它到 $X\times_Y X$ 的两个拉回相同, 故逆 $\beta_X$ 下降为 $\Hom(W,F)$ 在 $\id:Y\to Y$ 上的截面 $\beta$. 由于 $\Hom(F,F)$ 与 $\Hom(W,W)$ 都是层, 可得 $\beta\circ\alpha=\id_F$ 与 $\alpha\circ\beta=\id_W$. 另见~\cite[Tag 02W5]{SP}.
\end{proof}
